\documentclass[11pt]{article}
\usepackage[T1]{fontenc}
\usepackage[margin=0.82in]{geometry}
\usepackage{amsmath,amssymb,amsthm}
\usepackage{graphicx}
\usepackage{booktabs,tabularx}
\usepackage[hidelinks]{hyperref}
\usepackage{microtype}

\newtheorem*{theorem}{Counterexample theorem}
\newcommand{\R}{\mathbb R}
\newcommand{\F}{\mathcal F}
\newcommand{\Q}{\mathcal Q}

\title{A lacunary counterexample to the reverse weak difference inequality}
\author{New counterexample for the question in arXiv:2102.09657}
\date{13 August 2026}

\begin{document}
\maketitle

\begin{theorem}
Fix $N\ge1$, $1<p<\infty$, and $0<s<1$.  There is no constant $C$ such
that every
\[
 u\in\F^{-1}\!\left[C_c^\infty(\R^N\setminus\{0\})\right]
\]
satisfies
\begin{equation}\label{reverse}
 \|u\|_{\dot F^s_{p,2}}
 \le C\left[
   \frac{u(x)-u(y)}{|x-y|^{N/p+s}}
 \right]_{L^{p,\infty}(\R^N\times\R^N)}.
\end{equation}
Indeed, there are $u_m$ for which the norm on the left is
$\gtrsim\sqrt m$, while the weak norm on the right stays bounded.
Consequently the holomorphic-family formulation (1.29) in the source is
also false.
\end{theorem}

\section*{The source question}
Theorem~1.3 of Brezis--Seeger--Van Schaftingen--Yung proves the forward
bound
\[
 \left[\frac{u(x)-u(y)}{|x-y|^{N/p+s}}\right]_{L^{p,\infty}}
 \lesssim \|u\|_{\dot F^s_{p,2}}.
\]
The paragraph after it asks whether this can be reversed.  Using
$\dot F^s_{p,2}=[L^p,\dot W^{1,p}]_s$, it equivalently asks for a
holomorphic interpolation family whose boundary norms are bounded by the
weak difference quotient.

\begin{center}
\includegraphics[width=0.94\textwidth]{evidence/source_question_crop.png}

{\footnotesize Source evidence: arXiv:2102.09657, PDF page 7 (unaltered crop).}
\end{center}

\clearpage
\section*{Construction}
Choose a power $R=2^K\ge8$ and a nonzero Schwartz function $\phi$ such that
$\widehat\phi\in C_c^\infty(B(0,\varepsilon))$, with $\varepsilon>0$ small.
For $m\ge1$, put
\begin{equation}\label{um}
 u_m(x)=\phi(x)P_m(x_1),\qquad
 P_m(t)=\sum_{j=1}^m R^{-js}e^{2\pi iR^jt}.
\end{equation}
The Fourier support of $u_m$ is contained in the finite disjoint union
\[
 \bigcup_{j=1}^m B(R^je_1,\varepsilon),
\]
which is compact and avoids the origin.  Thus every $u_m$ belongs to the
source's test class.

\section*{The Triebel--Lizorkin norm diverges}
Use an equivalent smooth dyadic resolution adapted to the widely separated
frequency balls.  Its projection at level $Kj$ equals the identity on
$B(R^je_1,\varepsilon)$ and annihilates the other selected balls.  Hence
\[
 \Delta_{Kj}u_m=R^{-js}\phi(x)e^{2\pi iR^jx_1}.
\]
Keeping only these $m$ terms in the square function gives
\begin{align*}
 \|u_m\|_{\dot F^s_{p,2}}
 &\gtrsim
 \left\|\left(\sum_{j=1}^m
   \left|2^{Kjs}R^{-js}\phi(x)e^{2\pi iR^jx_1}\right|^2
 \right)^{1/2}\right\|_{L^p_x}\\
 &=\sqrt m\,\|\phi\|_{L^p},
\end{align*}
because $2^K=R$.  Equivalence of admissible Littlewood--Paley resolutions
has constants independent of the finite number $m$ of bands.

\section*{A uniform H\texorpdfstring{$\ddot{o}$}{o}lder bound}
The elementary geometric-series estimate
\begin{equation}\label{geom}
 \sum_{j=1}^m R^{-js}\min\{2,2\pi R^jr\}
 \le A_{R,s}r^s\qquad(r>0)
\end{equation}
holds uniformly in $m$.  To see it, split where $R^jr$ crosses $1$.
Before that scale, sum the geometric progression
$rR^{j(1-s)}$; after it, sum $R^{-js}$.  Both pieces are bounded by a
constant times $r^s$.  It follows that
\begin{equation}\label{holderP}
 |P_m(a)-P_m(b)|\le A_{R,s}|a-b|^s,
 \qquad |P_m(a)|\le (R^s-1)^{-1},
\end{equation}
uniformly in $m$.

Fix $M>N/p$ and let $w(x)=(1+|x|)^{-M}$.  Schwartz decay and the mean-value
theorem imply
\begin{equation}\label{phi}
 |\phi(x)|\le B_Mw(x),\qquad
 |\phi(x)-\phi(y)|
 \le B_M|x-y|^s\bigl(w(x)+w(y)\bigr).
\end{equation}
For the second bound, use the gradient estimate when $|x-y|\le1$ (points
on the segment have comparable weights) and the zeroth-order estimate when
$|x-y|>1$.

Writing $r=|x-y|$ and expanding the product in \eqref{um}, equations
\eqref{holderP}--\eqref{phi} yield
\begin{align}
 |u_m(x)-u_m(y)|
 &\le |\phi(x)-\phi(y)|\,|P_m(x_1)|
      +|\phi(y)|\,|P_m(x_1)-P_m(y_1)|\notag\\
 &\le C_{\phi,R,s,M}r^s\bigl(w(x)+w(y)\bigr),\label{pointwise}
\end{align}
again with no dependence on $m$.

\section*{The weak quotient stays bounded}
Divide \eqref{pointwise} by $r^{N/p+s}$.  It remains to estimate
\[
 H(x,y)=\frac{w(x)+w(y)}{|x-y|^{N/p}}.
\]
This weak estimate is explicit.  For every $\lambda>0$,
\begin{align*}
 \mathcal L^{2N}\!\left\{(x,y):
       \frac{w(x)}{|x-y|^{N/p}}>\lambda\right\}
 &=\int_{\R^N}\left|B\!\left(x,
       (w(x)/\lambda)^{p/N}\right)\right|dx\\
 &=\kappa_N\lambda^{-p}\|w\|_{L^p}^p.
\end{align*}
Applying this once to $w(x)$ and once to $w(y)$, with threshold
$\lambda/2$, gives
\[
 [H]_{L^{p,\infty}(\R^{2N})}
 \le 2(2\kappa_N)^{1/p}\|w\|_{L^p}<\infty.
\]
Therefore \eqref{pointwise} proves
\begin{equation}\label{uniformQ}
 \sup_{m\ge1}
 \left[\frac{u_m(x)-u_m(y)}{|x-y|^{N/p+s}}
 \right]_{L^{p,\infty}(\R^{2N})}<\infty.
\end{equation}
Together with the $\sqrt m$ lower bound, \eqref{uniformQ} contradicts every
putative constant in \eqref{reverse}.

\section*{Why the holomorphic-family version also fails}
If the family requested in source display (1.29) existed with a uniform
implicit constant, the definition of complex interpolation would give
\[
 \|u\|_{[L^p,\dot W^{1,p}]_s}
 \lesssim
 \left[\frac{u(x)-u(y)}{|x-y|^{N/p+s}}
 \right]_{L^{p,\infty}}.
\]
On the source's frequency-localized class,
$[L^p,\dot W^{1,p}]_s=\dot F^s_{p,2}$ with equivalent norms.  Taking
$u=u_m$ contradicts \eqref{uniformQ}.  Thus the reformulation fails along
with the direct reverse inequality.

\clearpage
\section*{Proof intuition}
Each frequency band is scaled so that it contributes one unit to the
Littlewood--Paley square function; $m$ orthogonal scales therefore accumulate
as $\sqrt m$.  The weak quotient behaves differently.  At a separation
$|x-y|\simeq R^{-k}$, the $k$th frequency is the largest relevant term;
lower and higher frequencies form convergent geometric tails.  Taking a
weak supremum sees the worst scale but does not add the $m$ scale
contributions.  This square-function versus supremum mismatch is exactly the
obstruction to reversal.

\paragraph{Scope and novelty.}
The counterexample works for every fixed $N\ge1$, $1<p<\infty$, and
$0<s<1$, using complex-valued Schwartz functions with compact Fourier
support away from zero.  It disproves the unrestricted reverse inequality
and the stated holomorphic-family bound.  Bounded exact-phrase and concept
searches through 13 August 2026 found no later resolution of this exact
question; human literature review is still recommended.

\paragraph{Human review.}
Check the adapted Littlewood--Paley projection on the separated balls, the
two geometric sums in \eqref{geom}, and the distribution calculation for
$H$.  These are the only substantive estimates.

\begin{thebibliography}{9}
\raggedright
\bibitem{source}
H. Brezis, A. Seeger, J. Van Schaftingen, and P.-L. Yung,
\emph{A new formula for the $L^p$ norm}, arXiv:2102.09657 (2021),
question following Theorem~1.3 on PDF page 7.
\end{thebibliography}

\end{document}
